\centerline{\bf \Huge Вписанные четырёхугольники.}

\begin{flushright}
        \scriptsize{Окружность — душа геометрии. Познайте окружность,\\ и вы не только познаете геометрию, но и возвысите душу свою…\\И.Ф. Шарыгин}
\end{flushright}


\textbf{Теорема.} Выпуклый четырёхугольник $ABCD$ можно вписать в окружность тогда и только тогда, когда выполнено хотя бы одно из эквивалентных условий:
\[
(1) \ \angle CAD = \angle CBD; \quad
(2) \ \angle ABD = \angle ACD; \quad
(3) \ \angle ABC + \angle ADC = 180^{\circ}.
\]

\begin{center}
\includegraphics[width=1\textwidth]{Летние Сборы/Сборы Элиста 2024/РВ/Geometry_9/Cyclic_qudrilateral.png}
\end{center}

\problem{1} 
На боковых сторонах $AB$ и $CD$ трапеции $ABCD$ выбраны точки $P$ и $Q$ соответственно. Оказалось, что четырехугольник $APQD$ вписанный. Докажите, что четырехугольник $BPQC$ вписанный.

\problem{2}
Внутри треугольника $ABC$ отмечена точка $P$, а на сторонах $BC, CA, AB$ — точки $A_1, B_1, C_1$ соответственно. Докажите, что если два из трех четырехугольников $AB_1PC_1, BC_1PA_1, CA_1PB_1$ вписанные, то третий тоже вписанный.

\problem{3}
Точка \( K \) — середина «меньшей» дуги \( AB \) вписанного четырёхугольника \( ABCD \). Хорды \( CK \) и \( AB \) пересекаются в точке \( P \), хорды \( DK \) и \( AB \) пересекаются в точке \( Q \). Доказать, что четырёхугольник \( CPQD \) вписанный.

\problem{4}
Точка \( O \) — центр описанной окружности равнобокой трапеции \( ABCD \) (\( AD \parallel BC \)), \( K \) — точка пересечения её диагоналей. Докажите, что точки \( A \), \( B \), \( K \), \( O \) лежат на одной окружности.

\problem{5}
На гипотенузе \( AB \) прямоугольного треугольника \( ABC \) во внешнюю сторону построен квадрат с центром в точке \( O \). Докажите, что \( CO \) — биссектриса прямого угла.

\problem{6}
В остроугольном треугольнике \( ABC \) с углом \( \angle A = 60^\circ \) отмечены точки \( H, I, O \) — ортоцентр и центры вписанной и описанной окружностей соответственно. Докажите, что описанная окружность треугольника \( HIO \) проходит через вершину \( B \).

\newpage

\hspace{3mm}

\centerline{\bf \Huge Вписанные четырёхугольники.}

\begin{flushright}
        \scriptsize{А теперь...моя очередь}
\end{flushright}

\problem{7}
\textbf{Лемма Фусса.} Окружности \( \omega_1 \) и \( \omega_2 \) пересекаются в точках \( A \) и \( B \). На \( \omega_1 \) выбираются точки \( X \neq A \) и \( Y \neq B \). Прямые \( XA \) и \( YB \) пересекают окружность \( \omega_2 \) в точках \( X_1 \) и \( Y_1 \) соответственно. Докажите, что \( XY \parallel X_1 Y_1 \).

\problem{8}
\textbf{Прямая Симсона.} Дан остроугольный треугольник \( ABC \). Точка \( P \) выбирается на меньшей дуге \( BC \) описанной около него окружности. Докажите, что проекции точки \( P \) на стороны треугольника лежат на одной прямой.

\problem{9}
\textbf{Лемма о куринной лапке.} Пусть у треугольника \( ABC \) точка \( I \) — центр вписанной окружности, точка \( I_a \) — центр вневписанной окружности, противоположной вершине \( A \), а точка \( L \) — точка пересечения отрезка \( II_a \) с дугой описанной окружности. Тогда точка \( L \) равноудалена от \( I \), \( I_a \), \( B \) и \( C \).

\problem{$10^{**}$}
\textbf{Люти гроб.} Дан остроугольный треугольник $ABC$. Окружность, проходящая через вершину $B$ и центр $O$ его описанной окружности, вторично пересекает стороны $BC$ и $BA$ в точках $P$ и $Q$ соответственно. Докажите, что точка пересечения высот треугольника $POQ$ лежит на прямой $AC$.